\section{Homogeneous limits and their fixed points}\label{sec:homogeneous}

The estimates of Section~\ref{sec:arithmetic} give probability limits and
place a restriction on their local concentration. We first use this
restriction to obtain positive entropy in almost every ergodic component.
Measure classification then reduces the problem to excluding proper block
orbits. The useful feature of those orbits is concrete: every point on one
is fixed by a nonidentity diagonal matrix with algebraic entries.

\subsection{Positive entropy in every positive-weight collection of components}

For an invariant probability $\nu$ and an invertible measure-preserving map
$T$, write $h_\nu(T)$ for its measure-theoretic entropy. Fix
$b\in L$ with distinct coordinates and put
\[
 \lambda=\min_{j\ne k}|b_j-b_k|>0,\qquad T(x)=xa(b).
\]
The next statement is the ordinary-ball argument of
\cite[Proposition~6.1 and Lemma~6.2]{PrimeDegreePackets}, included here
with its proof. It explains why a ball estimate for a measure can control
its ergodic components without passing that estimate to each individual
component.

\begin{proposition}\label{prop:positive-entropy}
Let $\mu$ be an $A$-invariant probability on $X$. Suppose that, for every
compact $\Omega\subset X$, there are $C_\Omega,r_\Omega>0$ such that
\begin{equation}\label{eq:ordinary-ball-hypothesis}
 \mu(xB(r))\le C_\Omega r^n
 \quad(x\in\Omega,\ 0<r<r_\Omega),
\end{equation}
where $B(r)$ is a matrix neighborhood of the identity of radius $r$.
Every $A$-invariant probability $\xi\le c\mu$, with $c<\infty$, satisfies
$h_\xi(T)\ge\lambda/3$. Consequently almost every $A$-ergodic component
of $\mu$ has positive entropy for $T$.
\end{proposition}

\begin{proof}
Choose a sufficiently small fixed $\epsilon>0$ and consider the two-sided
tube
\[
 \mathcal T_t=a(tb)B(\epsilon)a(-tb)
               \cap a(-tb)B(\epsilon)a(tb).
\]
If $h\in\mathcal T_t$ and $r=e^{-\lambda t}$, then
\[
 |h_{jj}-1|<\epsilon,\qquad |h_{jk}|\le\epsilon r\quad(j\ne k).
\]
Every nondiagonal term in the determinant expansion contains at least two
off-diagonal entries. Hence $\prod_jh_{jj}=1+O(r^2)$. The matrix
\[
 d(h)=\diag\left(h_{11},\ldots,h_{n-1,n-1},
                       \Bigl(\prod_{j<n}h_{jj}\Bigr)^{-1}\right)
\]
belongs to a fixed compact set $\mathcal D\subset A$ and satisfies
$d(h)^{-1}h\in B(Cr)$. An $r$-mesh of the $(n-1)$-dimensional set
$\mathcal D$ therefore covers $\mathcal T_t$ by at most
$Cr^{-(n-1)}$ sets $d_jB(C'r)$, with $d_j\in\mathcal D$.
For $x\in\Omega$, all centers $xd_j$ belong to the compact set
$\Omega\mathcal D$. Domination and
\eqref{eq:ordinary-ball-hypothesis} give
\begin{equation}\label{eq:limit-tube-bound}
 \xi(x\mathcal T_t)\le C_{\Omega,c}e^{-\lambda t}.
\end{equation}

We apply the entropy criterion of Einsiedler--Lindenstrauss--Michel--
Venkatesh~\cite[Corollary~3.3]{ELMV2009}. In the form needed here, it says
that if probabilities $\nu_i$ invariant under $a(\R b)$ converge to a probability
$\nu$, and $t_i\to\infty$, then
\[
 \nu_i\bigl(x[a(t_ib)Oa(-t_ib)\cap a(-t_ib)Oa(t_ib)]\bigr)
       \le C_\Omega e^{-2\eta t_i}\quad(x\in\Omega)
\]
for every compact $\Omega$, with an identity neighborhood $O$ allowed to
depend on $\Omega$, implies $h_\nu(T)\ge\eta$.
This criterion includes the noncompact space $X$ and requires that the
limit have total mass one. Take the constant sequence $\nu_i=\xi$,
$t_i=i$, and $\eta=\lambda/3$. Equation~\eqref{eq:limit-tube-bound}
verifies its hypothesis.

Write the $A$-ergodic decomposition as
$\mu=\int_Z\nu_z\,\dd\tau(z)$. The function
$z\mapsto h_{\nu_z}(T)$ is measurable: for increasing finite Borel
partitions $\mathcal P_k$ generating the Borel sets, it is given by
\[
 h_{\nu_z}(T)=\sup_k\inf_{N\ge1}\frac1N
 H_{\nu_z}\left(\bigvee_{j=0}^{N-1}T^{-j}\mathcal P_k\right).
\]
The entropy integral formula~\cite[Theorem~9.8]{Rokhlin1967} applies to
this decomposition into $T$-invariant probabilities. The component partition is
fixed by $T$, so the transformation on its quotient is the identity and
has zero entropy. The formula remains valid after restricting to a
measurable collection of components; the components need not themselves
be $T$-ergodic.

If the zero-entropy components had weight $w>0$, their normalized average
$\xi$ would satisfy $\xi\le w^{-1}\mu$ and $h_\xi(T)=0$. This
contradicts the first part of the proposition.
\end{proof}

\subsection{The algebraic alternatives}

We use the following measure-classification theorem in its full-degree
form.

\begin{theorem}[Einsiedler--Katok--Lindenstrauss]\label{thm:ekl}
For $n\ge3$, an $A$-invariant, $A$-ergodic probability on
$\SL_n(\Z)\backslash\SL_n(\R)$ that has positive entropy for some
diagonal one-parameter subgroup is the invariant probability on a closed
finite-volume orbit $xJ$, where $J$ is a closed connected subgroup
containing $A$.
\end{theorem}

This is \cite[Theorem~1.3]{EKL2006}, with the definition of an algebraic
measure given there. Their left-action convention is transferred to ours
by $\Gamma g\mapsto g^{-1}\Gamma$; this replaces $a(b)$ by $a(-b)$,
which has the same entropy. The Haar conclusion for prime $n$ in
\cite[Corollary~1.4]{EKL2006} does not identify the sextic alternatives.
We shall distinguish them using their fixed points.

\begin{lemma}\label{lem:block-groups}
Let $J$ be a connected closed subgroup of $G$ containing $A$, and suppose
$xJ$ carries a finite nonzero $J$-invariant measure. There is a partition
$\{1,\ldots,n\}=B_1\sqcup\cdots\sqcup B_d$ such that $J$ is the
identity component of the determinant-one block-diagonal group for this
partition.
\end{lemma}

\begin{proof}
The Lie algebra $\mathfrak j$ contains $L$ and is invariant under
$\Ad(A)$. Since the off-diagonal weight spaces are one-dimensional,
$\mathfrak j$ is the direct sum of $L$ and a collection of root lines
$\R E_{jk}$. Draw the edge $j\to k$ when this line occurs. The bracket
$[E_{ij},E_{jk}]=E_{ik}$ makes the graph transitively closed whenever
$i\ne k$.

A Lie group with a lattice is unimodular. Applying the vanishing of its
modular character to $A$ shows that
\[
 \sum_{j\to k}(b_j-b_k)=0\qquad(b\in L).
\]
Thus each vertex has equal indegree and outdegree. Every strongly
connected class is complete by transitivity. The graph of these classes
is acyclic. If it had an edge, a source class of an edge-containing
component would have strictly more outgoing than incoming edges in the
original graph, contradicting the sum of the degree balances over that
class. Hence there are no edges between classes.
The resulting Lie algebra is precisely the block-diagonal one. A
connected Lie subgroup is determined by its Lie algebra, giving the
assertion.
\end{proof}

The next lemma supplies a countable obstruction to each proper orbit.
It uses finite volume twice: first to identify a rational commutant and
then to exclude its idempotents.

\begin{lemma}\label{lem:fixed-pieces}
Every proper orbit $xJ$ in Theorem~\ref{thm:ekl} is contained in
$\Fix(a)=\{y\in X:ya=y\}$ for some nonidentity $a\in A$ whose entries
are algebraic numbers. There are only countably many possible matrices
$a$. For each such $a$, its fixed set is covered by countably many
smooth pieces. Whenever $a_{jj}\ne a_{kk}$, the direction of
$U_{jk}=\{1+tE_{jk}:t\in\R\}$ is nowhere tangent to these pieces,
and each piece meets a sufficiently small $U_{jk}$ plaque in at most
one point.
\end{lemma}

\begin{proof}
Write $x=\Gamma g$ and
$\Delta=g^{-1}\Gamma g\cap J$. This is a lattice in $J$.
For the rational linear algebra in this proof, use the auxiliary column
lattice $\Lambda=g^{-1}\Z^n$, which $\Delta$ preserves, and its rational
span $\Lambda_{\Q}$. This convention is separate from the row-lattice
realization of the point $x$.
We first show that the real matrix commutant of $\Delta$ equals the
commutant of $J$.
For any matrix $C$ commuting with $\Delta$, the map
\[
 \Delta\backslash J\longrightarrow\Mat_n(\R),\qquad
 \Delta j\longmapsto j^{-1}Cj
\]
is well-defined. It sends normalized Haar measure to a probability
invariant under conjugation by $A$. Choose a regular element of $A$.
Every off-diagonal matrix coordinate has a nonzero weight under its
conjugation. Recurrence of an invariant probability forces all such
coordinates to vanish almost surely: a nonzero coordinate expands
geometrically in one of the two time directions and cannot recur to a
bounded set. The map is continuous and Haar measure has full support,
so $j^{-1}Cj$ is diagonal for every $j\in J$.

Taking $j=1$ shows that $C$ is diagonal. Conjugating by a root subgroup
within one of the blocks of Lemma~\ref{lem:block-groups} forces the two
corresponding diagonal entries of $C$ to agree. Consequently the
commutant of $\Delta$ is exactly the algebra of block scalars, namely
the commutant of $J$.

Define
\[
 F=\{C\in\End_{\Q}(\Lambda_{\Q}):C\delta=\delta C
                                     \text{ for all }\delta\in\Delta\}.
\]
In the lattice coordinates supplied by $g$, these commutation relations
are rational linear equations, since $g\Delta g^{-1}\subset\Gamma$.
Only finitely many of these equations are needed in the finite-dimensional
matrix space. Thus $F$ has real extension
isomorphic to $\R^d$, where $d$ is the number of blocks. In particular
$F$ is a commutative, reduced, finite-dimensional $\Q$-algebra, a product
of totally real fields.

Suppose $F$ had a nontrivial idempotent. Its image in real coordinates
would select a nonempty proper union of blocks and hence a subspace $V$
defined over the rational lattice. The group $\Delta$ preserves the
full lattice in $V$ and its inverse does too. Thus
\[
 \chi(j)=|\det(j|_V)|
\]
is a character of $J$ trivial on $\Delta$. It is nontrivial on $A$,
because both the selected and the unselected union of blocks are
nonempty. It is therefore onto $\R_{>0}$. Pushing the finite invariant
measure on $\Delta\backslash J$ forward by $\chi$ would give a finite
nonzero multiplicatively invariant measure on $\R_{>0}$, an
impossibility. Hence $F$ is a field.

The representation of $F$ on the rational lattice space has some
dimension $m$ over $F$. After extending scalars to $\R$, each embedding
of $F$ has multiplicity $m$, so all blocks have size $m$ and $dm=n$.
The elements of $F$ carrying $\Lambda$ into itself form an order.
Indeed, in lattice coordinates this multiplier set is represented by
$(gFg^{-1})\cap\Mat_n(\Z)$, hence is discrete. Clearing the denominators of a
rational basis gives a full-rank submodule. Thus the multiplier set is
a finitely generated group of rank $d$, and it is a subring containing $1$.
Dirichlet's unit theorem for this order gives unit rank $d-1$.

If $J\ne G$, then $d>1$. Choose a unit of infinite order and take an even
power so that every real embedding is positive. Its action on the
lattice is an automorphism with determinant
$N_{F/\Q}(u)^m=1$. In real coordinates it is a matrix $a\in A$,
constant on each block and therefore central in $J$. It fixes $x$,
and centrality gives $(xj)a=(xa)j=xj$ for every $j\in J$.
Its entries are conjugates of an algebraic unit. The collection of
possible such diagonal matrices is countable.

Finally, fixing $a$ and lifting to $G$, the fixed-point equation is
\[
 ga g^{-1}=\gamma\qquad\text{for some }\gamma\in\Gamma.
\]
For each $\gamma$, the nonempty solution set is a coset of the
centralizer of $a$. It is smooth, with tangent root directions exactly
those satisfying $a_{jj}=a_{kk}$. Restricting to countably many
injective local quotient charts gives the asserted pieces in $X$.
If $a_{jj}\ne a_{kk}$, two points of one local $U_{jk}$ plaque cannot
solve the same lifted equation: their relative element $1+tE_{jk}$
would commute with $a$, forcing $t=0$. This proves the last assertion.
\end{proof}

For $n=6$, the equality $dm=6$ leaves only block sizes $3+3$,
$2+2+2$, or $1+1+1+1+1+1$ for a proper finite-volume orbit.
The last case is $J=A$, whose translation action on a compact quotient
has zero entropy. Thus the positive-entropy alternatives that must be
excluded have two blocks of size three or three blocks of size two.

By Proposition~\ref{prop:ball}, Proposition~\ref{prop:positive-entropy},
and Theorem~\ref{thm:ekl}, almost every ergodic component of any packet
limit is one of these homogeneous probabilities. To complete the proof,
it will suffice to show that such a limit assigns no mass to the
countable fixed-point pieces in Lemma~\ref{lem:fixed-pieces}. We now
construct conditional measures that can detect those pieces while
retaining the arithmetic information at the growing times of
Proposition~\ref{prop:bowen}.
